% !TeX root = Report.tex
\chapter{Описание алгоритмов и схема оценивания}

В данной главе излагаются алгоритмы, используемые далее в численном исследовании. Основное внимание уделяется тем свойствам рассматриваемых методов, которые существенны для задачи оценивания параметров прямой методом Хафа после предварительного подавления шума.

\section{Модель изображения и параметризация прямой}

Рассматривается изображение в градациях серого, представленное матрицей
\[
\mathbf{M} = (M_{ij}) \in \mathbb{R}^{m \times n}.
\]
Предполагается, что наблюдаемое изображение раскладывается в сумму компоненты сигнала и шумовой компоненты:
\[
\mathbf{M} = \mathbf{S} + \mathbf{R},
\]
где $\mathbf{S}$ --- компонента сигнала, а элементы шумовой матрицы $\mathbf{R}$ независимы и имеют распределение $\mathcal{N}(0,\sigma^2)$. В данной работе предполагается, что пикселям матрицы $\mathbf{S}$, принадлежащим прямой, соответствует фиксированная положительная интенсивность $c_k$, а вне прямой значения равны нулю.

В основной постановке численного эксперимента рассматривается одна прямая, задаваемая в декартовых координатах уравнением
\[
y = ax + b.
\]
Такая запись естественна для постановки численного эксперимента, поскольку коэффициенты $a$ и $b$ непосредственно задают положение прямой на изображении. При описании преобразования Хафа далее используется параметризация
\[
\rho = x \cos \theta + y \sin \theta,
\]
где $\rho$ --- расстояние от начала координат до прямой, а $\theta$ --- угол нормали к прямой \cite{duda1972hough}.

\begin{figure}[H]
	\centering
	\includegraphics[width=0.5\textwidth]{rt_params.pdf}
	\caption{Параметризация прямой через $(\rho,\theta)$}
	\label{fig:rt-params}
\end{figure}

\section{Преобразование Хафа}

Преобразование Хафа --- классический метод оценки параметров прямых и других простых геометрических объектов по набору активных точек изображения \cite{duda1972hough}. Активными далее называются те пиксели, интенсивность которых после предобработки больше нуля. Идея метода состоит в переходе от исходного пространства изображения к пространству параметров прямой. Для каждой активной точки $(x_0,y_0)$ множество параметров всех прямых, проходящих через нее, задается кривой
\[
\rho(\theta) = x_0 \cos \theta + y_0 \sin \theta.
\]
Если несколько точек лежат на одной прямой, то соответствующие кривые в пространстве $(\rho,\theta)$ пересекаются в общей точке, которая и задает параметры этой прямой.

В рассматриваемой дискретной постановке пространство параметров заменяется сеткой с шагами $h_\rho$ и $h_\theta$:
\[
\rho_i = -\rho_{\max} + i h_\rho, \qquad \theta_j = j h_\theta.
\]
Далее строится массив $A=(A_{ij})$, где значение в ячейке $(i,j)$ показывает, сколько кривых вида $\rho(\theta)$ пересекаются вблизи точки $(\rho_i,\theta_j)$:
\[
A_{ij} =
\sum_{(x,y)\in E}
\mathbf{1}
\left\{
\left|
\rho_i - \bigl(x\cos\theta_j + y\sin\theta_j\bigr)
\right|
\le \frac{h_\rho}{2}
\right\}.
\]
Здесь
\[
E = \{(x,y)\colon \widetilde{M}_{xy} > 0\}
\]
обозначает множество активных пикселей после предобработки. Чем больше значение $A_{ij}$, тем более вероятно, что точки в исходном пространстве лежат на прямой с параметрами, близкими к $(\rho_i,\theta_j)$.

Таким образом, преобразование Хафа в рассматриваемой работе состоит из следующих шагов:
\begin{enumerate}
	\item к изображению применяется выбранный метод предобработки и формируется множество активных пикселей $E$;
	\item для каждой активной точки $(x,y)\in E$ строится кривая $\rho(\theta)=x\cos\theta + y\sin\theta$ в пространстве параметров, в соответствующую позицию массива $A$ добавляются единицы;
	\item по точке максимума массива $A$ находятся оценки параметров $(\hat{\rho},\hat{\theta})$.
\end{enumerate}

Если после шумоподавления в изображении сохраняется много ложных активных точек, то в пространстве параметров возникают дополнительные максимумы, и точность итоговых оценок прямой ухудшается. Поэтому качество предобработки непосредственно влияет на качество оценивания параметров.

\begin{figure}[H]
	\centering
	\includegraphics[width=0.82\textwidth]{ACCUMULATOR_pic.pdf}
	\caption{Пример аккумуляторного массива для изображения с несколькими прямыми}
	\label{fig:accumulator}
\end{figure}

\section{Базовый метод SSA и его комплекснозначный вариант}

Метод анализа сингулярного спектра SSA предназначен для представления ряда в виде суммы интерпретируемых компонент, например тренда, периодических составляющих и шума \cite{ts_structure,ssa_with_r}. Комплекснозначный вариант Complex SSA обобщает эту конструкцию на случай комплексных рядов и особенно удобен в задачах, где после преобразования Фурье сигнал естественно становится комплекснозначным \cite{zhornikova2016}.

Рассмотрим комплексный временной ряд
\[
z = (z_1,\dots,z_N), \qquad z_k \in \mathbb{C}.
\]
Алгоритм CSSA, как и классический SSA, состоит из четырех основных этапов.

\textbf{1. Вложение.}
Выбирается длина окна $L$, где $1 < L < N$, и строится траекторная матрица
\[
\mathbf{X} = \mathcal{T}_L(z) =
\begin{pmatrix}
	z_1 & z_2 & \dots & z_K \\
	z_2 & z_3 & \dots & z_{K+1} \\
	\vdots & \vdots & \ddots & \vdots \\
	z_L & z_{L+1} & \dots & z_N
\end{pmatrix},
\qquad K = N-L+1.
\]
Эта матрица ганкелева, то есть элементы на ее побочных диагоналях совпадают.

\textbf{2. Сингулярное разложение.}
Для траекторной матрицы строится матрица
\[
\mathbf{S} = \mathbf{X}\mathbf{X}^{*},
\]
где $\mathbf{X}^{*} = \overline{\mathbf{X}}^{T}$ --- эрмитово сопряженная матрица, где $\overline{A}$ обозначает комплексное сопряжение элементов произвольной матрицы $A$. Тогда S является эрмитовой и неотрицательно определенной. Поэтому все ее собственные числа неотрицательны и могут быть упорядочены по невозрастанию:
\[
\lambda_1 \ge \lambda_2 \ge \dots \ge \lambda_L \ge 0.
\]
Если $d=\max\{i\colon \lambda_i>0\}$, то сингулярное разложение записывается в виде
\[
\mathbf{X} = \sum_{i=1}^{d}\mathbf{X}_i
=
\sum_{i=1}^{d}\sqrt{\lambda_i}\,U_iV_i^{*},
\]
где $U_i$ и $V_i$ --- левые и правые сингулярные векторы. Числа $\sqrt{\lambda_i}$ называются сингулярными числами. Наборы $(\sqrt{\lambda_i},U_i,V_i)$ принято называть собственными тройками.

\textbf{3. Группировка.}
Выделяется множество индексов $I$, соответствующее компонентам сигнала, и строится сумма
\[
\mathbf{X}_I = \sum_{i\in I}\mathbf{X}_i.
\]
Именно на этом шаге отделяется предполагаемая часть ряда, соответствующая сигналу, от шумовой.

\textbf{4. Восстановление.}
К матрице $\mathbf{X}_I$ применяется диагональное усреднение, в результате чего получается восстановленный ряд
\[
\widehat{z} = \mathcal{T}_L^{-1}(\mathbf{X}_I).
\]

\subsection*{Свойства CSSA, используемые далее}

В дальнейшем для интерпретации результатов будут использоваться следующие стандартные свойства CSSA \cite{ts_structure,ssa_with_r,zhornikova2016}, а также на некоторые свойства конкретных временных рядов.

\begin{enumerate}
	\item Все собственные числа матрицы $\mathbf{S}=\mathbf{X}\mathbf{X}^{*}$ неотрицательны, а собственные тройки можно упорядочить по убыванию вклада в разложение.
	\item Если наблюдаемый ряд представим в виде суммы компоненты сигнала и шумовой компоненты, то при не слишком большом шуме наибольшие сингулярные числа обычно связаны с сигналом, а маленькие --- преимущественно с шумом. Именно этим мотивируется восстановление по небольшому числу компонент, упорядоченных по убыванию.
	\item Число ненулевых сингулярных чисел совпадает с рангом траекторной матрицы. Следовательно, ранг задает минимальное число элементарных компонент, необходимых для точного представления сигнала.
	\item Для невырожденных гармонических и экспоненциально-модулированных сигналов траекторная матрица имеет малый ранг (1 или 2). В вещественном SSA гармоническая компонента обычно порождает пару близких компонент, тогда как в Complex SSA в некоторых конфигурациях комплексная гармоника может иметь ранг 1. Это свойство принципиально важно для рассмотрения одной прямой в следующей главе.
	\item Качество восстановления определяется не только величиной собственных чисел, но и корректной группировкой компонент. Поэтому в задачах шумоподавления недостаточно механически увеличивать число сохраняемых компонент: это может возвращать вместе с сигналом и большое количество шума, а в некоторых случаях возвращать только шум.
	\item Если отделимость сигнала от шума нарушается, компоненты начинают смешиваться, что приводит к ''растеканию'' сигнала и к появлению ложных пикселей после восстановления.
\end{enumerate}

Последние два свойства играют существенную роль в дальнейшем изложении. Во второй главе они используются при обсуждении выбора одной компоненты для \texttt{CSSA row.row}, а в третьей --- при интерпретации причин, по которым визуально качественное восстановление не всегда приводит к повышению точности оценок параметров прямой методом Хафа.
% Не очень нравится абзац выше, TODO: исправить.

Для рассматриваемой задачи существенным является следующее соображение: если сигнал в частотной области представим суммой небольшого числа комплексных экспонент, то его траекторная матрица обладает малым эффективным рангом. Поэтому восстановление по небольшому числу компонент оказывается естественным механизмом шумоподавления. На этом основано применение CSSA к строкам и столбцам изображения после преобразования Фурье. В случае одной прямой в ряде ситуаций можно ожидать, что основная структура сигнала восстанавливается по одной компоненте; данный вопрос подробно рассматривается в следующей главе.

\begin{figure}[H]
	\centering
	\includegraphics[width=0.72\textwidth]{CSSA_pic.pdf}
	\caption{Иллюстрация обработки изображения методом на основе CSSA}
	\label{fig:cssa-scheme}
\end{figure}

\section{Четыре варианта CSSA для изображений}

В работе рассматривается матрица изображения $\mathbf{M}\in\mathbb{R}^{m\times n}$ как система одномерных рядов. Предобработка на основе CSSA использует прямое преобразование Фурье по строкам, затем применяет CSSA либо к строкам, либо к столбцам полученной комплексной матрицы, а после этого выполняет обратное преобразование Фурье по одному из измерений. Именно выбор измерения для CSSA и обратного преобразования и порождает четыре варианта алгоритма.

Обозначим через $\mathcal{F}_{\mathrm{row}}$ прямое преобразование Фурье по строкам матрицы, через $\mathcal{F}^{-1}_{\mathrm{row}}$ и $\mathcal{F}^{-1}_{\mathrm{col}}$ --- обратные преобразования по строкам и столбцам, а через $\mathcal{C}_{\mathrm{row}}^{(r)}$ и $\mathcal{C}_{\mathrm{col}}^{(r)}$ --- применение CSSA к каждой строке или каждому столбцу с сохранением $r$ компонент.

Тогда рассматриваются следующие четыре варианта:
\begin{align*}
\widehat{\mathbf{S}}_{\mathrm{row.row}}
&=
\operatorname{Re}\!\Bigl(
\mathcal{F}^{-1}_{\mathrm{row}}
\bigl(
\mathcal{C}_{\mathrm{row}}^{(r)}(\mathcal{F}_{\mathrm{row}}(\mathbf{M}))
\bigr)
\Bigr), \\
\widehat{\mathbf{S}}_{\mathrm{row.col}}
&=
\operatorname{Re}\!\Bigl(
\mathcal{F}^{-1}_{\mathrm{col}}
\bigl(
\mathcal{C}_{\mathrm{row}}^{(r)}(\mathcal{F}_{\mathrm{row}}(\mathbf{M}))
\bigr)
\Bigr), \\
\widehat{\mathbf{S}}_{\mathrm{col.row}}
&=
\operatorname{Re}\!\Bigl(
\mathcal{F}^{-1}_{\mathrm{row}}
\bigl(
\mathcal{C}_{\mathrm{col}}^{(r)}(\mathcal{F}_{\mathrm{row}}(\mathbf{M}))
\bigr)
\Bigr), \\
\widehat{\mathbf{S}}_{\mathrm{col.col}}
&=
\operatorname{Re}\!\Bigl(
\mathcal{F}^{-1}_{\mathrm{col}}
\bigl(
\mathcal{C}_{\mathrm{col}}^{(r)}(\mathcal{F}_{\mathrm{row}}(\mathbf{M}))
\bigr)
\Bigr).
\end{align*}
% TODO: написать более красиво, как в презентации

Обозначение \texttt{row.row} означает, что после прямого преобразования Фурье по строкам метод CSSA также применяется к строкам, а обратное преобразование Фурье выполняется по строкам. Аналогично интерпретируются комбинации \texttt{row.col}, \texttt{col.row} и \texttt{col.col}.

Такое построение связано с тем, что после перехода в частотную область строка изображения, содержащая один яркий пиксель, близка к комплексной экспоненте, а сумма небольшого числа пикселей приводит к сумме нескольких экспонент \cite{trickett2003,zhornikova2016}. Следовательно, применение CSSA после преобразования Фурье имеет теоретическое основание как метод выделения низкоранговой структуры сигнала и подавления шума.

С точки зрения дальнейшего анализа важны два обстоятельства. Во-первых, разные варианты CSSA используют одно и то же исходное изображение, но по-разному организуют восстановление структуры после частотного преобразования. Во-вторых, даже при одинаковом числе сохраняемых компонент они могут существенно отличаться по тому, какие остаточные пиксели остаются после предобработки. Именно это затем влияет на аккумуляторный массив метода Хафа и на точность оценивания параметров прямой.

\begin{algorithm}[H]
	\caption{Алгоритм \texttt{CSSA ROW-ROW}}
	\begin{algorithmic}[1]
		\Require матрица изображения $\mathbf{M}\in\mathbb{R}^{m\times n}$, число сохраняемых компонент $r$
		\Ensure восстановленное изображение $\widehat{\mathbf{S}}$
		\For{$p=1,\dots,m$}
			\State вычислить ДПФ строки: $\boldsymbol{\xi}^{(p)} \gets \mathcal{F}(\mathbf{m}^{(p)})$
			\State применить CSSA к ряду $\boldsymbol{\xi}^{(p)}$
			\State сохранить $r$ компонент и получить восстановленный ряд $\widehat{\boldsymbol{\xi}}^{(p)}$
			\State вычислить обратное ДПФ:
			$\widehat{\mathbf{s}}^{(p)} \gets \operatorname{Re}(\mathcal{F}^{-1}(\widehat{\boldsymbol{\xi}}^{(p)}))$
		\EndFor
		\State составить матрицу $\widehat{\mathbf{S}}$ из строк $\widehat{\mathbf{s}}^{(1)},\dots,\widehat{\mathbf{s}}^{(m)}$
		\State \Return $\widehat{\mathbf{S}}$
	\end{algorithmic}
\end{algorithm}

\section{Классические методы предобработки}

Помимо алгоритмов на основе CSSA, в работе рассматриваются стандартные методы шумоподавления, применяемые к той же зашумленной матрице изображения перед использованием метода Хафа.

\subsection*{Медианный фильтр}

Медианный фильтр заменяет значение каждого пикселя медианой значений в окне \cite{huang1979}. Если через $W_{ij}$ обозначить квадратное окно вокруг пикселя $(i,j)$ размера $k$, то результат фильтрации имеет вид
\[
\widehat{M}_{ij} = \operatorname{median}\{M_{uv} \colon (u,v)\in W_{ij}\}.
\]
Если не оговорено иное, используется окно размера \(3\times 3\).
Существенное свойство медианного фильтра состоит в способности подавлять одиночные выбросы, не размывая резкие границы столь же сильно, как линейные сглаживающие фильтры. Далее рассматривается стандартный вариант с квадратным локальным окном.

\subsection*{Фильтр Винера}

Фильтр Винера в общем виде строит линейную оценку сигнала, минимизирующую среднеквадратичную ошибку \cite{wiener1949}. Далее рассматривается локально-адаптивная модификация этого подхода для изображений; соответствующее описание приводится, например, у Lim \cite{lim1990}. Термин ``локальный'' означает, что параметры фильтра, как и в случае медианного фильтра, определяются отдельно для каждого пикселя по его окну.

Пусть $W_{ij}$ --- квадратное окно размера $k\times k$ с центром в точке $(i,j)$. Стандартным считается окно размера \(5\times 5\). Тогда для каждого пикселя сначала вычисляются локальное среднее
\[
\mu_{ij} = \frac{1}{|W_{ij}|}\sum_{(u,v)\in W_{ij}} M_{uv}
\]
и локальная дисперсия
\[
\nu_{ij} =
\frac{1}{|W_{ij}|}\sum_{(u,v)\in W_{ij}} M_{uv}^2 - \mu_{ij}^2.
\]
После этого вводится величина $\sigma_n^2$, обозначающая оценку дисперсии шумовой компоненты изображения. В рассматриваемой схеме эта величина определяется формулой
\[
\sigma_n^2 = \frac{1}{mn}\sum_{i=1}^{m}\sum_{j=1}^{n} \nu_{ij},
\]
то есть как среднее значение локальных дисперсий по всем пикселям изображения. Далее результат записывается как
\[
\widehat{M}_{ij}
=
\mu_{ij}
+
\frac{\max(\nu_{ij}-\sigma_n^2,\,0)}{\max(\nu_{ij},\,\varepsilon)}
\left(M_{ij}-\mu_{ij}\right).
\]
Здесь $\varepsilon>0$ --- малая константа, предотвращающая деление на нуль. Приведенная формула показывает, что при малой локальной дисперсии, сравнимой с оценкой шума, поправочный множитель оказывается малым, и значение пикселя смещается к локальному среднему. Если же локальная дисперсия велика, степень сглаживания уменьшается, что способствует сохранению локальной структуры изображения.

Таким образом, под фильтром Винера далее понимается локально-адаптивный фильтр, а не глобальный частотный вариант и не задача винеровской деконволюции. Такой фильтр согласуется с локальной структурой изображения: в почти однородных областях он подавляет шум сильнее, а в областях с заметной вариативностью в большей степени сохраняет полезный сигнал.

\subsection*{BM3D}

Метод BM3D относится к числу наиболее известных алгоритмов подавления гауссовского шума на изображениях \cite{dabov2007}. Его основная идея состоит в использовании самоподобия изображения: небольшие фрагменты, расположенные в различных областях, нередко обладают близкой структурой. Это позволяет обрабатывать не отдельные фрагменты, а группы взаимно похожих фрагментов.

Пусть выбран фрагмент изображения фиксированного размера, называемый опорным фрагментом. Для него в некоторой области поиска отыскиваются другие фрагменты, близкие к нему по выбранной мере сходства (например, MSE или SSIM). Близость фрагментов означает, что их значения пикселей и соответствующая геометрическая структура, например граница, линия или однородная область, согласованы с точностью до допустимого уровня искажений.

Найденные фрагменты объединяются в трехмерный массив: первые две координаты задают положение пикселя внутри фрагмента, а третья координата нумерует отобранные фрагменты. Дальнейшая обработка состоит из двух основных стадий:
% Далее к такой группе применяется совместное преобразование, при котором полезный сигнал приобретает более компактное представление, чем шум.
\begin{enumerate}
	\item на первой стадии выполняется жесткая пороговая фильтрация коэффициентов преобразования, что позволяет подавить шумовые составляющие;
	\item на второй стадии выполняется винеровская фильтрация коэффициентов с использованием оценки, полученной на первом этапе.
\end{enumerate}

После обратного преобразования получают несколько оценок одних и тех же пикселей, поскольку различные группы похожих фрагментов перекрываются. Эти оценки объединяются в итоговое изображение посредством взвешенного усреднения. Тем самым формируется единая восстановленная матрица изображения.

Благодаря использованию самоподобия различных областей изображения метод BM3D способен подавлять шум эффективнее чисто локальных методов, особенно при наличии повторяющихся структур.

\subsection*{Квантильное отсечение по порогу}

В качестве простого метода отсечения по порогу рассматривается также вариант, в котором порог задается квантилем. Пусть $Q_{0.9}(\mathbf{M})$ обозначает $0.9$-квантиль яркостей изображения. Тогда результат определяется формулой
\[
\widehat{M}_{ij} =
\begin{cases}
M_{ij}, & M_{ij} \ge Q_{0.9}(\mathbf{M}), \\
0, & M_{ij} < Q_{0.9}(\mathbf{M}).
\end{cases}
\]
В отличие от предыдущих методов, этот способ не сглаживает изображение, а лишь сохраняет наиболее яркие пиксели. Он используется как простая схема отсечения по порогу, поскольку позволяет оценить, в какой степени один лишь глобальный порог способен улучшить работу метода Хафа.

Для всех рассматриваемых методов принципиально важно, что после предобработки изображение далее интерпретируется через множество активных пикселей. Поэтому для метода Хафа существенна не столько сама интенсивность пикселя, сколько факт его участия в дальнейшем голосовании.

\section{Схема оценивания}

Во всех экспериментах используется одна и та же схема оценивания:
\begin{enumerate}
	\item к зашумленному изображению применяется один из методов предобработки;
	\item по обработанному изображению формируется множество активных пикселей;
	\item по этому множеству строится аккумуляторный массив метода Хафа;
	\item по максимумам аккумулятора находятся оценки $(\hat{\rho},\hat{\theta})$.
\end{enumerate}

Ключевой момент состоит в том, что завершающий этап оценивания фиксирован, а изменяется только предобработка. Поэтому различия в точности итоговых оценок могут быть интерпретированы как следствие различий между методами шумоподавления.
